\section{Resolved bridge records and open independent strengthenings}

The central construction, intrinsic nearby-cycle/Leray proof, explicit cusp
collapse, based van Kampen calculation, sphere recognition,
algebraic-dimension calculation, completed local Kummer map,
cone-compactified sphere carriers, moving four-sphere slices, fixed-target
metric scope, fixed-energy counting law, full-torus magnetic operator, and
finite compact-link Gaussian system are now verified in this workbench.  The
following list preserves C1--C3, MV1, and L1 as resolved audit records and
adds the exact K1, C4--C10, S1--S2, and G1--G2 records.  Each closure uses
actual primitive maps; matching ranks
alone was not accepted.  The global cover of \(X\) remains absent, the
literal fixed-target cusp extension remains \Refuted, and the physical
Yang--Mills bridge remains \Open.

\subsection{Period orientation and the modular deck certificate}

\begin{auditrecord}[G0: simultaneous finite signs and exact deck data]
\Verified\ Let $g_0=(g_1g_2)^{-1}$.  The two finite generators are
normalized simultaneously, not independently, by
\[
 g_1'(z_1)=e^{-2\pi i/3},\qquad
 g_2'(z_2)=e^{-2\pi i/4}=-i,
 \qquad g_0\text{ parabolic}.
\]
For
\[
 S_1(w)=\frac{w-1}{w},\qquad S(w)=-\frac1w,
 \qquad P(w)=w-1,
\]
the normalized period lift has the exact deck certificate
\[
 \varphi(g_1)=S_1,\qquad \varphi(g_2)=S,\qquad
 \varphi(g_0)=P,\qquad S_1SP=1
 \quad\text{in }\operatorname{PSL}_2(\mathbb Z).
\]
The quadratic ramification at $z_2$ proves only
$(g_2'(z_2))^2=-1$; the sign $-i$ is forced by retaining the chosen
$g_1$ orientation and the parabolic cusp class simultaneously.  See
\cref{lem:triangle-generator-orientations,eq:period-deck-certificate}.
\end{auditrecord}

\subsection{Completed local Kummer cover}

\begin{auditrecord}[K1: exact Kummer extension through the toric cusp]
\Verified\ Fix \(N\geq1\), put \(t=v^N\), and do not reuse the target
action at the source.  The forced source factor and source action are
\[
 C_N(v)=\frac1N C(v^N),\qquad
 c^{(N)}_{\bar\lambda}(v)=
 \exp\bigl(2\pi iC_N(v)\bar\lambda\bigr),
\]
\[
 \Psi^{(N)}_{\bar\lambda}(\xi,v)=
 \bigl(c^{(N)}_{\bar\lambda}(v)
 v^{B_0\bar\lambda}\xi,v\bigr).
\]
After shrinking the \(v\)-disc, this is a free properly discontinuous
action on the same smooth toric model \(Y^{(N)}\), and
\[
 f_N:N_0^{\langle N\rangle}=Y^{(N)}/\bar\Lambda
 \longrightarrow\Delta_v
\]
 is proper with smooth total space.  Its central fibre
 \(W_N:=(f_N^{-1}(0))_{\mathrm{red}}\) is reduced,
irreducible, and nonnormal, with normalization the toric del Pezzo surface
of degree six.

Multiplication by \(N\) on the full three-dimensional toric lattice gives
\[
 P_N(\xi_1,\xi_2,v)=(\xi_1^N,\xi_2^N,v^N).
\]
The equality
\[
 P_N\Psi^{(N)}_{\bar\lambda}=\Psi_{\bar\lambda}P_N
\]
uses exactly

\[
 \bigl(c^{(N)}_{\bar\lambda}(v)\bigr)^N
 =c_{\bar\lambda}(v^N);
\]
without this modified factor the power map does not descend in general.
On a maximal toric chart the local-ring map is
\[
 \mathbb C\{z_0,z_1,z_2\}\longrightarrow
 \mathbb C\{w_0,w_1,w_2\},\qquad z_i\longmapsto w_i^N,
\]
free of rank \(N^3\).  Hence it descends to the finite flat map
\[
 \kappa_N:N_0^{\langle N\rangle}\longrightarrow N_0,\qquad
 f_0\circ\kappa_N=\rho_N\circ f_N,\qquad \deg\kappa_N=N^3,
\]
and the Jacobian calculation gives
\[
 \kappa_N^*W=NW_N,\qquad
 \operatorname{Ram}(\kappa_N)=(N-1)W_N.
\]

The base change \(N_0[N]=N_0\times_{\Delta_t}\Delta_v\) is normal, and
\(\kappa_N\) factors uniquely through
\[
 Q_N:N_0^{\langle N\rangle}\longrightarrow N_0[N].
\]
At a triple point its exact ring map is
\[
 \frac{\mathbb C\{z_0,z_1,z_2,v\}}
 {(z_0z_1z_2-v^N)}
 \longrightarrow\mathbb C\{w_0,w_1,w_2\},\qquad
 z_i\longmapsto w_i^N,\qquad v\longmapsto w_0w_1w_2.
\]
The source is the invariant ring for
\[
 G_N=\{(\varepsilon_0,\varepsilon_1,\varepsilon_2)\in\mu_N^3:
 \varepsilon_0\varepsilon_1\varepsilon_2=1\}\simeq\mu_N^2.
\]
Thus \(Q_N\) has generic degree \(N^2\).  For \(N>1\) it is quasi-etale,
but it is not flat: its non-etale and nonflat loci are exactly the inverse
images of the double and triple strata.  The decisive closed-fibre lengths
are
\[
 2N-1>N\quad\text{on a double stratum},
 \qquad
 3N^2-3N+1>N^2\quad\text{at a triple point}.
\]
It is therefore not a normalization map.

On normalized central fibres, \(Q_N\) lifts to the toric map
\([N]:\mathrm{dP}_6\to\mathrm{dP}_6\) of degree \(N^2\), ramified with
index \(N\) on each of the six boundary curves and compatible with the
paired conductors and both triple-point classes.  On a regular fibre it is
the sublattice cover
\[
 \mathbb Z\langle\widehat\gamma,\widehat u,
 N\widehat w,N\widehat\delta\rangle\subset\Lambda,
 \]
with deck group \((\mathbb Z/N\mathbb Z)^2\) and
\[
 S_N^{-1}M_0^NS_N=M_0,\qquad S_N=\operatorname{diag}(1,1,N,N).
\]
This closes the extension across the \emph{local} toric cusp.  It does not
construct a cover over the entire base or a cover of \(X\), and it supplies
no compatible completion through the two finite fillings.  See
\cref{cusp-cover-extension-theorem,cusp-cover-naive-action-failure,%
cusp-cover-global-scope}.
\end{auditrecord}

\subsection{Cone-compactified sphere carriers}

\begin{auditrecord}[C4: exact \(S^{12}\) and \(S^{24}\) smash routes]
\Verified\ For
\[
 K_\Lambda=\mathbb Z\langle\widehat w,\widehat\delta\rangle,
 \qquad \overline T^2=
 \frac{(\Lambda/K_\Lambda)\otimes\mathbb R}{\Lambda/K_\Lambda},
\]
the identity of \(M_0\) on \(\Lambda/K_\Lambda\) gives a real-analytic
fibrewise group epimorphism
\[
 \overline\pi:N_0\setminus W\longrightarrow D_0^*\times\overline T^2
\]
with connected \(T^2\)-kernel.  On the radial fibre it is exactly the native
quotient \(\pi_T:F_T\to Q_T\), and
\(\operatorname{sys}(Q_T)\geq T-C\).

Put
\[
 \mathcal A=
 \frac{[0,\epsilon]\times S^1\times\overline T^2}
 {\{0\}\times S^1\times\overline T^2\sim v_0,
  \ \{\epsilon\}\times S^1\times\overline T^2\sim v_1}
 \cong\operatorname{Susp}(T^3).
\]
The formulas
\[
 q_2(x)=[r,e^{i\phi},\overline\pi_f(x)],\qquad
 q_2(W)=v_0,\qquad
 q_2(X\setminus\operatorname{int}N_0)=v_1
\]
define a continuous surjective based quotient \(q_2:X\to\mathcal A\).
Its smash square is
\[
 q_2^{\wedge2}:X\wedge X\cong S^{12}\longrightarrow
 \mathcal A\wedge\mathcal A,
\]
and on \(F_T^2\) it is exactly
\(\pi_T\times\pi_T:F_T^2\to Q_T^2\).  In the intrinsic flat quotient
metric, \(Q_T^2\to\mathbb R^4\) pointed locally isometrically, smoothly,
and in pointed Gromov--Hausdorff topology.

The full-monodromy-invariant character gives the localized target
\[
 \mathcal B=\operatorname{Susp}(T^2)
\]
and a continuous surjective based quotient
\[
 q_{1,\mathrm{cusp}}:X\longrightarrow\mathcal B.
\]
Its fourth smash power has source \(S^{24}\) and restricts on synchronized
radial fibres to
\[
 \chi_{\gamma,T}^{\times4}:F_T^4\longrightarrow C_T^4.
\]

For the full-base version, glue the pair-of-pants circle target to the two
finite Seifert targets:
\[
 \mathcal M=(P\times S^1)
 \underset{h_1^{S^1}}\cup C_1
 \underset{h_2^{S^1}}\cup C_2,\qquad
 \partial\mathcal M\cong T^2,
\]
and change the cusp codomain by putting
\[
 \mathcal C=\mathcal M\cup_{\partial\mathcal M}
 \operatorname{Cone}(\partial\mathcal M).
\]
The regular-core, finite-filling, and cusp formulas glue to a continuous
surjective based quotient \(q_1:X\to\mathcal C\), with \(W\) sent to the
cone vertex.  Its fourth smash power is
\[
 q_1^{\wedge4}:X^{\wedge4}\cong S^{24}\longrightarrow
 \mathcal C^{\wedge4},
\]
and on the radial fibres it is exactly
\(\chi_{\gamma,T}^{\times4}:F_T^4\to C_T^4\), where
\[
 H_{\mathrm{cusp}}:=\operatorname{Im}h(0),\qquad
 B_{\mathrm{cusp}}:=\operatorname{Im}b(0),
 \qquad
 \operatorname{sys}(C_T^4)=\ell_T
 =T+H_{\mathrm{cusp}}-B_{\mathrm{cusp}}+O(T^{-1}),
 \qquad (C_T^4,0)\longrightarrow(\mathbb R^4,0)
\]
in all three pointed intrinsic topologies above.

The target homology is exact:
\[
 \widetilde H_k(\mathcal A)=
 \begin{cases}\mathbb Z^3&k=2,3,\\\mathbb Z&k=4,\\0&\text{otherwise},\end{cases}
\]
\[
 \widetilde H_k(\mathcal A\wedge\mathcal A)=
 \begin{cases}
 \mathbb Z^9&k=4,\\\mathbb Z^{18}&k=5,\\
 \mathbb Z^{15}&k=6,\\\mathbb Z^6&k=7,\\\mathbb Z&k=8,\\
 0&\text{otherwise},
 \end{cases}
\]
\[
 \widetilde H_k(\mathcal B)=
 \begin{cases}\mathbb Z^2&k=2,\\\mathbb Z&k=3,\\0&\text{otherwise},\end{cases}
\]
\[
 \widetilde H_k(\mathcal B^{\wedge4})=
 \begin{cases}
 \mathbb Z^{16}&k=8,\\\mathbb Z^{32}&k=9,\\
 \mathbb Z^{24}&k=10,\\\mathbb Z^8&k=11,\\\mathbb Z&k=12,\\
 0&\text{otherwise},
 \end{cases}
\]
and
\[
 \widetilde H_k(\mathcal C)=
 \begin{cases}\mathbb Z&k=2,3,\\0&\text{otherwise},\end{cases}
 \qquad
\widetilde H_k(\mathcal C^{\wedge4})=
 \begin{cases}
 \mathbb Z&k=8,12,\\\mathbb Z^4&k=9,11,\\
 \mathbb Z^6&k=10,\\0&\text{otherwise}.
 \end{cases}
\]
The point-set quotient theorems and these homology calculations do not
determine the based homotopy classes of
\[
 q_2^{\wedge2}:S^{12}\longrightarrow\mathcal A\wedge\mathcal A,\qquad
 q_{1,\mathrm{cusp}}^{\wedge4}:S^{24}\longrightarrow\mathcal B^{\wedge4},
 \qquad
 q_1^{\wedge4}:S^{24}\longrightarrow\mathcal C^{\wedge4}.
\]
All three conical-target classes remain \Open.
Thus the targets are not \(S^4\), \(S^{12}\), \(S^{24}\), or compact
four-manifolds.  The source spheres are smash powers, not Cartesian powers.
The Euclidean limits are intrinsic to the varying flat slices; in every
fixed compatible metric on the compact targets, the slices collapse to the
cone vertex.  See
\cref{cusp-cone-S12-route,cusp-cone-S24-local-route,%
cusp-cone-S24-route,cusp-cone-metric-scope}.
\end{auditrecord}

\begin{auditrecord}[C5: fixed-target and finite-cover obstruction boundary]
\Verified\ For every primitive connected fibrewise torus quotient of \(r\)
copies over the full regular base or a finite cover, divergence of the cusp
systole forces
\[
 \operatorname{rank}(\Lambda^{\oplus r}/N)\leq r.
\]
Indeed bounded cusp periods force \(K_\Lambda^{\oplus r}\subset N\), and
the finite-index monodromy orbit
\[
 (A_2A_1)^n\widehat w
 =\widehat w+n(\widehat u-\widehat\delta)
\]
then forces \(\widehat u\) into each primitive kernel.  In particular, no
nontrivial finite base cover preserves \(K_\Lambda\), and two copies cannot
give a full-base expanding four-torus quotient.

The fixed-target cohomology obstructions remain exact: no map to a fixed
\(\mathbb T^4\) retains the four independent target characters, and no map to a
fixed \(S^4\) retains a degree-one collapse of \(Q_T^2\) or \(C_T^4\).
Moreover
\[
 X^{\times2}\simeq S^6\times S^6\not\simeq S^{12},\qquad
 X^{\times4}\simeq(S^6)^{\times4}\not\simeq S^{24},
\]
the synchronized fibre products have real dimensions \(10\) and \(18\),
and
\[
 X*X\cong S^{13},\qquad X^{*4}\cong S^{27}.
\]
All connected unbranched finite covers of the compact sphere sources are
trivial, and neither \(S^{12}\) nor \(S^{24}\) admits an almost-complex
structure.

Therefore the literal extension to the old fixed torus or sphere target
remains \Refuted.  The maps in C4 do not reverse that result: they replace
the target at the cusp and collapse the entire boundary torus to a new cone
 stratum.  No global base cover, cover of \(X\), or branched or toroidal
 global modification
or compact complex sphere has been constructed.  See
\cref{cusp-cone-obstructions}.
\end{auditrecord}

\begin{auditrecord}[C6: sphere-to-sphere quotients with moving metric
four-spheres]
\Verified\ Fix the reconstructed based diffeomorphism
\(X\cong_{\mathrm{diff}}S^6\).  Collapsing the coordinate one-skeleton in
\(\overline T^2\) and then reducing the suspension gives a continuous
surjective based quotient
\[
 \widehat q_3=\widehat p_2\circ q_2:(X,x_0)\longrightarrow(S^3,*).
\]
Forgetting the cusp-base angle in the localized invariant-character target
gives a second continuous surjective based quotient
\[
 \widehat q_\gamma=\widehat p_1\circ q_{1,\mathrm{cusp}}:
 (X,x_0)\longrightarrow(S^2,*).
\]
Their smash powers are therefore quotient maps
\[
 \widehat q_3^{\wedge2}:S^{12}\twoheadrightarrow S^6,
 \qquad
 \widehat q_\gamma^{\wedge4}:S^{24}\twoheadrightarrow S^8.
\]
Both full sphere-valued maps are based null-homotopic.  First write
\[
 g=\widehat q_3:S^6\longrightarrow S^3.
\]
Hatcher's full low-dimensional homotopy-group table
\cite[p.~339]{Hatcher2002} gives
\[
 \pi_6(S^3)=\mathbb Z/12,
 \qquad
 \pi_{12}(S^6)=\mathbb Z/2,
\]
so \(12g=0\).  Toda's Proposition~5.11
\cite[Proposition~5.11]{Toda1962} identifies the latter group as
\[
 \pi_{12}(S^6)=\mathbb Z/2\{\nu_6\circ\nu_9\},
\]
and shows that suspension carries this generator through the groups
\(\pi_{n+6}(S^n)=\mathbb Z/2\{\nu_n\circ\nu_{n+3}\}\), \(n\geq5\).
Consequently stabilization gives an isomorphism
\[
 \operatorname{st}:\pi_{12}(S^6)\xrightarrow{\ \cong\ }\pi_6^{\mathrm S}.
\]
Isaksen--Wang--Xu
\cite[Table~1, p.~113]{IsaksenWangXu2023} gives
\[
 \pi_3^{\mathrm S}=\mathbb Z/24\{\nu\},
 \qquad
 \pi_6^{\mathrm S}=\mathbb Z/2\{\nu^2\}.
\]
Because \(12\operatorname{st}(g)=0\), there is
\(a\in\mathbb Z/12\) for which
\(\operatorname{st}(g)=2a\nu\).  Multiplicativity of stabilization for
smash products therefore gives
\[
 \operatorname{st}(g^{\wedge2})
 =\operatorname{st}(g)^2
 =4a^2\nu^2=0.
\]
Injectivity of the displayed stabilization isomorphism proves
\[
 [g^{\wedge2}]=0\in\pi_{12}(S^6).
\]
For every \(r>2\), smashing this based null-homotopy with
\(g^{\wedge(r-2)}\) proves
\[
 [g^{\wedge r}]=0\in\pi_{6r}(S^{3r}).
\]

For the second map write
\[
 f=\widehat q_\gamma:S^6\longrightarrow S^2.
\]
Hatcher's stable-range result
\cite[Corollary~4.24, p.~360]{Hatcher2002}, applied with \(n=8\) and
\(k=4\), identifies
\(\pi_{12}(S^8)\) with the fourth stable stem
\(\pi_4^{\mathrm S}\).  Table~1 of Isaksen--Wang--Xu
\cite[p.~113]{IsaksenWangXu2023} gives the complete equality
\[
 \pi_{12}(S^8)\cong\pi_4^{\mathrm S}=0.
\]
Thus the based map
\[
 f\wedge\operatorname{id}_{S^6}=\Sigma^6f:
 S^6\wedge S^6\longrightarrow S^2\wedge S^6
\]
is null-homotopic.  Under the canonical based smash identifications,
\[
 f^{\wedge2}
 =\bigl(\operatorname{id}_{S^2}\wedge f\bigr)
  \circ\bigl(f\wedge\operatorname{id}_{S^6}\bigr),
\]
so the displayed null-homotopy of the right-hand factor gives a based
null-homotopy of \(f^{\wedge2}\).  For every \(r>2\), associativity gives
\[
 f^{\wedge r}=f^{\wedge2}\wedge f^{\wedge(r-2)},
\]
and the smash of a based null-homotopy with the based map
\(f^{\wedge(r-2)}\) is again a based null-homotopy.  Therefore
\[
 [f^{\wedge r}]=0\in\pi_{6r}(S^{2r})
 \qquad\text{for every }r\geq2.
\]
In particular both full continuous surjective carriers
\[
 \widehat q_3^{\wedge2}:S^{12}\twoheadrightarrow S^6,
 \qquad
 \widehat q_\gamma^{\wedge4}:S^{24}\twoheadrightarrow S^8
\]
are based null-homotopic.  Surjectivity and null-homotopy are compatible
and concern different properties of the same maps.  Neither calculation
determines the one-copy classes by itself. The separate explicit cap in
Theorem~\ref{cusp-null:thm} now proves \([\widehat q_3]=0\);
\([\widehat q_\gamma]\in\pi_6(S^2)\) remains open. Nor does either smash calculation
determine the based homotopy classes of the conical-target smash maps
\(q_2^{\wedge2}\), \(q_{1,\mathrm{cusp}}^{\wedge4}\), or
\(q_1^{\wedge4}\); those remain \Open.  In particular, null-homotopy after
a sphere-valued postcomposition cannot be pulled back through that
postcomposition.
On synchronized radial fibres the exact restrictions are
\[
 F_T^{\times2}\longrightarrow Q_T^{\times2}\longrightarrow
 S_T^2\wedge S_T^2=:\mathbb S_{12,T}^4
\]
and
\[
 F_T^{\times4}\xrightarrow{\chi_{\gamma,T}^{\times4}}
 C_T^{\times4}\longrightarrow
 C_T^{\wedge4}=:\mathbb S_{24,T}^4,
\]
where the second arrow collapses the corresponding fat wedge.  With the
quotient length metrics and midpoint basepoints, every radius-\(R\) ball in
\(\mathbb S_{12,T}^4\) is exactly Euclidean when \(T-C>4R\), and every
radius-\(R\) ball in \(\mathbb S_{24,T}^4\) is exactly Euclidean when
\(\ell_T>4R\).  Thus both moving four-spheres converge to \(\mathbb R^4\)
pointed locally isometrically, smoothly, and in pointed
Gromov--Hausdorff topology.  Sliding the reduced-suspension coordinate to
an endpoint gives based null-homotopies of both moving-sphere inclusions;
their fundamental classes die in the fixed targets.  Under the unscaled
Euclidean quotient metric, rank proves that two copies are minimal for the
native cusp-local quotient and four copies are minimal for a rank-four
quotient invariant over the full regular base.  See
\cref{cusp-cone-optimal-exponents,%
cusp-cone-angle-forgotten-sphere-carriers,%
cusp-cone-S12-sphere-map-null,cusp-cone-S24-sphere-map-null,%
cusp-cone-sphere-map-homotopy-status,%
cusp-cone-moving-S4-null-homotopy}.
\end{auditrecord}

\begin{auditrecord}[C7: fixed-target metric realization and its exact
dimensional boundary]
\Verified\ In every fixed metric compatible with a compact cone-target
topology, the synchronized radial slices converge in Hausdorff, hence
Gromov--Hausdorff, topology to the smash basepoint.  This topological
collapse is distinct from blowing up one fixed cusp-adapted metric.  With
\(r_T=e^{-2\pi T}\), the fixed tensors
\[
 g_{\mathcal A}=dr^2+r^2du^2+\frac{r^2}{T(r)^2}h_{T(r)},
 \qquad
 g_{\mathcal D}=dr^2+r^2(du^2+d\theta^2),
 \quad \mathcal D\in\{\mathcal B,\mathcal C\},
\]
extend to compatible length metrics on the compact targets.  Their induced
slice forms obey
\[
 \left(\frac{T}{r_T}\right)^2j_T^*g_{Z_{12}}
 =g_{Q_T}\oplus g_{Q_T},
 \qquad
 \left(\frac{\ell_T}{r_T}\right)^2
 (k_T^{\mathcal D})^*g_{Z_{24}^{\mathcal D}}
 =g_{C_T}^{\oplus4}.
\]
The corresponding slice bubbles are \(\mathbb R^4\), while the full
ambient bubbles under the same rescalings are \(\mathbb R^8\) for the
\(S^{12}\) route and \(\mathbb R^{12}\) for either \(S^{24}\) route.
For any scalar rescaling \(\lambda_T\), put
\(\beta_T=\lambda_Tr_T\).  The complete slice trichotomy is a point if
\(\beta_T\to0\), a square flat four-torus of period \(L\) if
\(\beta_T\to L\in(0,\infty)\), and \(\mathbb R^4\) if
\(\beta_T\to\infty\).  The unit-volume limit is the square flat four-torus,
not \(\mathbb R^4\); every rescaling producing \(\mathbb R^4\) sends the
cone vertex to infinite distance.  Hence the verified \(\mathbb R^4\) is a
moving slice bubble, neither a vertex tangent cone nor a compact-target
asymptotic cone, and no metric on a source sphere is being asserted.  See
\cref{sphere-metric-uniform-collapse,sphere-metric-fixed-realization,%
sphere-metric-scaling-trichotomy,sphere-metric-unit-volume,%
sphere-metric-tangent-scope}.
\end{auditrecord}

\begin{auditrecord}[C8: the finite-cover \(\mathbb R^4\) theorem is exact
but ancillary]
\Verified\ For \(N\geq1\), put
\[
 S_N=\operatorname{diag}(1,1,N,N),\qquad
 \widetilde F_{T,N}=\mathbb R^4/P_TS_N\mathbb Z^4.
\]
The identity on \(\mathbb R^4\) induces an \(N^2\)-sheeted holomorphic
local-isometric cover \(\widetilde F_{T,N}\to F_T\) with deck group
\((\mathbb Z/N\mathbb Z)^2\).  After the cyclic cusp-base change
\(t_c=s^N\),
\[
 S_N^{-1}M_0^NS_N=M_0,
\]
so the covers assemble over the punctured changed base.  Their exact lattice
estimate is
\[
 \operatorname{sys}(\widetilde F_{T,N})\geq\min\{T-C,N\};
\]
hence every \(T_j,N_j\to\infty\) gives a pointed locally isometric, smooth,
and Gromov--Hausdorff limit \(\mathbb R^4\).  Its character lattice has
momenta \(P_T^{-\mathsf T}S_N^{-1}k\), so modes not pulled back from
\(F_T\) fill the two previously discrete momentum directions.

This theorem is not the requested sphere-carrier generalization of the
native \(\mathbb T^2\times\mathbb R^2\) quotient: it changes the fibre
lattice and makes a cyclic cusp-base change, and it constructs neither an
\(S^{12}\) nor an \(S^{24}\) carrier.  It is retained as a correct
ancillary result and is not substituted for C6--C7.  Its local completion is
K1; neither statement constructs a global cover of \(X\).
\end{auditrecord}

\begin{auditrecord}[C9: the boundary-family comparison is not a quotient]
\Refuted For \(\nu\in\{12,24\}\), the constructed map
\[
 \beta_\nu:\mathfrak E_\nu\longrightarrow\mathcal D_\nu
\]
is continuous and surjective, restricts to a homeomorphism
\[
 \mathfrak E_\nu\setminus p_\nu^{-1}(0)
 \xrightarrow{\ \cong\ }
 \mathcal D_\nu\setminus\{d_\nu\},
\]
and satisfies
\(\beta_\nu^{-1}(d_\nu)=p_\nu^{-1}(0)\cong\mathbb R^4\).
It is nevertheless nonproper and is not a quotient map.  Indeed, choose
\(a\in\mathbb R^4\setminus\mathbb Z^4\), let \(s_j\downarrow0\), and set
\[
 z_j=[s_j,A_{1/s_j}^{(\nu)}a]\in\mathfrak E_\nu.
\]
For \(\nu=12\),
\[
 \inf_{n\in\mathbb Z^4}
 \|A_{1/s_j}^{(12)}(a+n)\|
 \geq(1/s_j-C)\operatorname{dist}(a,\mathbb Z^4)\longrightarrow\infty,
\]
and for \(\nu=24\) the same infimum is
\[
 \ell_{1/s_j}\operatorname{dist}(a,\mathbb Z^4)\longrightarrow\infty.
\]
Thus \(Z=\{z_j:j\geq1\}\) is closed in \(\mathfrak E_\nu\).  It is
\(\beta_\nu\)-saturated, because \(\beta_\nu\) is injective off the zero
fibre and \(\beta_\nu(Z)\) omits \(d_\nu\).  But
\[
 \beta_\nu(z_j)=[e^{-2\pi/s_j},a]\longrightarrow d_\nu,
\]
so \(\beta_\nu(Z)\) is not closed in \(\mathcal D_\nu\).  A quotient map
would make a set closed exactly when its inverse image is closed; hence
\(\beta_\nu\) cannot be quotient.  Consequently
\(\mathcal D_\nu\not\cong
\mathfrak E_\nu/p_\nu^{-1}(0)\) through \(\beta_\nu\), and the map is not a
topological blow-down or resolution.

The positive theorem is narrower and unchanged: \(\mathfrak E_\nu\) is a
Hausdorff second-countable topological five-manifold with boundary
\(p_\nu^{-1}(0)\cong\mathbb R^4\); its positive fibres are the native
expanding four-tori and converge to that boundary fibre pointed locally
isometrically, smoothly, and in Gromov--Hausdorff topology.  The moving
four-sphere results in C6 and the fixed-metric scope in C7 do not use the
refuted quotient interpretation.  See
\cref{cusp-cone-exact-decompactification-space}.
\end{auditrecord}

\begin{auditrecord}[C10: literal native \(\mathbb R^4\) open cells in the
sphere targets]
\Verified Put
\[
 U_{12,T}:=\mathbb S_{12,T}^4\setminus\{*\},
 \qquad
 U_{24,T}:=\mathbb S_{24,T}^4\setminus\{*\}.
\]
If a quotient collapses a closed subset \(K\) to one point, its restriction
\[
 X\setminus K\longrightarrow (X/K)\setminus\{*\}
\]
is a homeomorphism: the source is saturated and open, the restricted map is
quotient and bijective, and every bijective quotient map is a homeomorphism.
Applying this first to the transported coordinate one-skeleton gives
\[
 Q_T\setminus K_T^{(1)}
 \cong S_T^2\setminus\{*\}
 \cong (0,1)^2\cong\mathbb R^2.
\]
Likewise \(C_T\setminus\{0\}\) is an open interval and hence is
homeomorphic to \(\mathbb R\).  Because the complement of the smash
basepoint is exactly the product of the complements of the factor
basepoints, the quotient-stage restrictions are homeomorphisms
\[
 (Q_T\setminus K_T^{(1)})^2
 \xrightarrow{\ \cong\ }U_{12,T}\cong\mathbb R^4,
\]
\[
 (C_T\setminus\{0\})^4
 \xrightarrow{\ \cong\ }U_{24,T}\cong\mathbb R^4.
\]
The moving four-spheres are topologically embedded, so composition gives
topological embeddings
\[
 \mathbb R^4\cong U_{12,T}\hookrightarrow S^6,
 \qquad
 \mathbb R^4\cong U_{24,T}\hookrightarrow S^8.
\]

These are nonclosed open cells.  Each is open in its moving four-sphere and
has that whole four-sphere as its closure, because the omitted smash
basepoint is a limit point.  Neither is open in its ambient sphere:
otherwise it would be simultaneously a four-manifold and an open six- or
eight-manifold, contradicting invariance of dimension.  Nor is either cell
globally isometric to Euclidean \(\mathbb R^4\) in the assigned compact
quotient metric or the restriction of a fixed compatible compact-target
metric, since those metrics are bounded while Euclidean \(\mathbb R^4\) is
unbounded.  The exact metric content is the midpoint-ball exhaustion:
for every fixed \(R\), once \(T-C>4R\), respectively \(\ell_T>4R\), the
corresponding radius-\(R\) pointed ball is exactly and onto isometric to the
Euclidean radius-\(R\) ball.

This result uses the native quotient stages.  It neither invokes the
ancillary finite-cover unspooling nor depends on the nonquotient comparison
map \(\beta_\nu\) of C9.  See
\cref{cusp-cone-literal-R4-open-cells,%
cusp-cone-S12-onto-exact-ball,cusp-cone-S24-onto-exact-ball}.
\end{auditrecord}

\subsection{Fixed-energy and full-torus magnetic spectra}

\begin{auditrecord}[S1: exact fixed-energy scalar count]
\Verified\ Put
\[
 \mathcal D_T:=|\det P_T|=\operatorname{Vol}(F_T^{\mathrm{flat}}),
 \qquad
 N_T(E):=\#\{k\in\mathbb Z^4:
 4\pi^2\lVert P_T^{-\mathsf T}k\rVert^2\leq E\}.
\]
For every fixed \(E\geq0\), exact elimination of the two integral
coordinates fixed by the constant periods reduces the other two coordinates
to translated ellipses and gives
\[
 N_T(E)=\pi\mathcal D_T
 \sum_{n\in\mathbb Z^2}
 \left(\frac{E}{4\pi^2}-\lVert n\rVert^2\right)_+
 +O_E(T).
\]
The sum is finite.  When \(0<E<4\pi^2\), only \(n=0\) contributes, so
\[
 N_T(E)=\frac{E\mathcal D_T}{4\pi}+O_E(T)
       =\frac{ET^2}{4\pi}+O_E(T).
\]
The error follows from the planar area--perimeter lattice estimate and
\(\lVert\mathcal A_T\rVert=O(T)\).  This is a marked flat-fibre count, not
a Weyl law for a fixed metric on the compact total space.  See
\texttt{s6\_spectral\_degeneration\_preprint.tex},
Proposition~\texttt{prop:fixed-energy-count}.
\end{auditrecord}

\begin{auditrecord}[S2: exact full-torus integral magnetic spectrum]
\Verified\ For every integral alternating form
\(\mathsf E\in\bigwedge^2(\mathbb Z^4)^*\), the factor of automorphy
\[
 J_n^{\mathsf E}(u)=
 \exp\!\left(-2\pi i\sum_{j<k}E_{jk}n_ju_k\right)
\]
defines a Hermitian line bundle \(\mathcal L_{\mathsf E,T}\), and
\[
 \nabla^{\mathsf E}=d+2\pi i\sum_{j<k}E_{jk}u_j\,du_k,
 \qquad F_{\nabla^{\mathsf E}}=2\pi i\mathsf E.
\]
Its magnetic form is closed on \(H^1\); the Friedrichs operator has domain
\(H^2\), compact resolvent, and, for nondegenerate \(\mathsf E\), spectrum
\[
 \operatorname{Spec}(\Delta_{\mathsf E,T})
 =\{(2n_1+1)b_1+(2n_2+1)b_2:n_1,n_2\geq0\},
\]
where \(\{\pm ib_1,\pm ib_2\}\) is the spectrum of
\(G_T^{-1/2}(2\pi\mathsf E)G_T^{-1/2}\).  Retaining every determinant
factor gives
\[
 b_1b_2=\frac{4\pi^2|\operatorname{Pf}(\mathsf E)|}{\mathcal D_T},
 \qquad
 \operatorname{gap}(\Delta_{\mathsf E,T})
 \leq4\pi\sqrt{\frac{|\operatorname{Pf}(\mathsf E)|}{\mathcal D_T}}
 =O(T^{-1}).
\]
This construction is on the full marked fibre and does not use the
exceptional first-coordinate elliptic quotient.  It does not assert
full-monodromy invariance of \(\mathsf E\), a global line bundle on \(X\),
or a Yang--Mills mass-gap theorem.  See
\texttt{s6\_spectral\_degeneration\_preprint.tex},
Theorem~\texttt{thm:full-torus-magnetic-spectrum} and
Corollary~\texttt{cor:full-torus-magnetic-gap}.
\end{auditrecord}

\subsection{Compact links and the finite Gaussian boundary}

\begin{auditrecord}[G1: exact cubulation, holonomy, and periodic flux]
\Verified\ Keep \(M_{\mathrm{cov}}\geq1\) and
\(N_{\mathrm{mesh}}\geq2\) independent, and put
\[
 V_{N_{\mathrm{mesh}}}=(\mathbb Z/N_{\mathrm{mesh}}\mathbb Z)^4,
 \qquad
 E_{N_{\mathrm{mesh}}}=V_{N_{\mathrm{mesh}}}\times\{1,2,3,4\}.
\]
The map
\[
 \Phi_{T,M_{\mathrm{cov}}}:\mathbb R^4/\mathbb Z^4
 \longrightarrow\widetilde F_{T,M_{\mathrm{cov}}},\qquad
 [y]\longmapsto[P_{T,M_{\mathrm{cov}}}y]
\]
is a diffeomorphism.  It transports the standard \(N_{\mathrm{mesh}}\)-cubical
subdivision to the whole cover torus, and on vertices gives the exact
bijection
\[
 (\mathbb Z/N_{\mathrm{mesh}}\mathbb Z)^4
 \xrightarrow{\ \cong\ }
 \widetilde F_{T,M_{\mathrm{cov}}}[N_{\mathrm{mesh}}],\qquad
 [n]\longmapsto
 \left[\frac1{N_{\mathrm{mesh}}}P_{T,M_{\mathrm{cov}}}n\right].
\]
For an edge \(\gamma_{n,j}\), parallel transport with
\[
 U'=-A(\dot\gamma_{n,j})U,\qquad U(0)=I_2
\]
defines
\[
 \operatorname{Hol}_{T,M_{\mathrm{cov}},N_{\mathrm{mesh}}}:
 \Omega^1(\widetilde F_{T,M_{\mathrm{cov}}};\mathfrak g)
 \longrightarrow G^{E_{N_{\mathrm{mesh}}}},
\]
and for \(A^g=g^{-1}Ag+g^{-1}dg\) its exact gauge law is
\[
 \operatorname{Hol}(A^g)_{n,j}
 =g(\gamma_{n,j}(1))^{-1}\operatorname{Hol}(A)_{n,j}
  g(\gamma_{n,j}(0)).
\]

With \(H_{\mathrm{link}}=\operatorname{diag}(i,-i)\) and
\(N=N_{\mathrm{mesh}}\), the explicit
periodic inverse-transport links satisfy
\[
 U_{12}(n)=\exp(2\pi H_{\mathrm{link}}/N^2)
 \quad\text{for every }n,
 \qquad U_{ij}(n)=I_2\quad((i,j)\ne(1,2),\ i<j).
\]
The principal logarithm is unambiguous and gives
\[
 \mathcal F_N^{\log}(n)=2\pi H_{\mathrm{link}}\,dy^1\wedge dy^2,
 \qquad
 S^{\log}_{T,M_{\mathrm{cov}},N}(U)
 =\frac{2\pi^2M_{\mathrm{cov}}^2}{g_{\mathrm{YM}}^2D_T}.
\]
The exact metric minor is
\[
 (G_{T,M_{\mathrm{cov}}}^{-1})_{11}
 (G_{T,M_{\mathrm{cov}}}^{-1})_{22}
 -(G_{T,M_{\mathrm{cov}}}^{-1})_{12}^2=D_T^{-2},
\]
and the explicitly declared \((1,2)\)-sector Wilson functional is
\[
 S^{W,12}_{T,M_{\mathrm{cov}},N}(U)
 =\frac{M_{\mathrm{cov}}^2N^4}{2g_{\mathrm{YM}}^2D_T}
 \left(2-2\cos\frac{2\pi}{N^2}\right),
\]
with limit \(2\pi^2M_{\mathrm{cov}}^2/(g_{\mathrm{YM}}^2D_T)\).
This does not assert a canonical scalar Wilson action for every
nonorthogonal metric.  See
\cref{cusp-link-cubical-realization,cusp-link-holonomy-map,%
cusp-link-flux-theorem}.
\end{auditrecord}

\begin{auditrecord}[G2: finite Gaussian oscillator and exact physical boundary]
\Verified\ On the reduced Coulomb space
\[
 \mathscr C_{T,M_{\mathrm{cov}},N_{\mathrm{mesh}}}
 =\ker d_{N_{\mathrm{mesh}}}^*\cap
 (\ker\Delta^{(1)}_{N_{\mathrm{mesh}}})^\perp,
\]
the operator
\[
 \mathcal L_{T,M_{\mathrm{cov}},N_{\mathrm{mesh}}}
 =g_{\mathrm{YM}}^{-2}d_{N_{\mathrm{mesh}}}^*
 d_{N_{\mathrm{mesh}}}
\]
with componentwise Fourier symbol
\[
 \mathfrak d_{N_{\mathrm{mesh}}}(k)_r
 =N_{\mathrm{mesh}}
 \left(e^{2\pi i k_r/N_{\mathrm{mesh}}}-1\right)
\]
has exact positive Fourier spectrum
\[
 \lambda(k)=\frac1{g_{\mathrm{YM}}^2}
 \left\lVert P_T^{-\mathsf T}S_{M_{\mathrm{cov}}}^{-1}
 \mathfrak d_{N_{\mathrm{mesh}}}(k)\right\rVert^2,\qquad k\ne0,
\]
with complex transverse multiplicity \(9\).  In particular,
\[
 \lambda(\mathbf e_1)=
 \frac{4N_{\mathrm{mesh}}^2\sin^2(\pi/N_{\mathrm{mesh}})}
 {g_{\mathrm{YM}}^2}
 \frac{L_T^2+m_T^2}{D_T^2},
\]
independently of \(M_{\mathrm{cov}}\).

The Friedrichs realization of
\[
 \widehat H^{\mathrm{Gauss}}
 =-\tfrac12\Delta_{\mathscr C_{T,M_{\mathrm{cov}},N_{\mathrm{mesh}}}}
 +\tfrac12\langle
 a,\mathcal L_{T,M_{\mathrm{cov}},N_{\mathrm{mesh}}}a\rangle
\]
has a unique Gaussian ground state and
\[
 \operatorname{gap}(\widehat H^{\mathrm{Gauss}})
 =\sqrt{\lambda_{\min}(\mathcal L)}.
\]
For the exact diagonal sequence
\[
 T_j^{\mathrm{rad}}=j^2,\qquad M_{\mathrm{cov},j}=j,\qquad
 N_{\mathrm{mesh},j}=j^3,
\]
one has simultaneously
\[
 \operatorname{sys}(\widetilde F_{T_j^{\mathrm{rad}},M_{\mathrm{cov},j}})
 \to\infty,
 \qquad
 \max_r\frac{\lVert P_{T_j^{\mathrm{rad}},M_{\mathrm{cov},j}}\mathbf e_r\rVert}
 {N_{\mathrm{mesh},j}}\to0,
\]
\[
 \operatorname{gap}(\widehat H^{\mathrm{Gauss}})\to0,
 \qquad
 S^{\log}_{T_j^{\mathrm{rad}},M_{\mathrm{cov},j},N_{\mathrm{mesh},j}}(U)
 \to0.
\]
Thus the cover geometry tends pointed locally isometrically to
\(\mathbb R^4\) while its finite mesh resolves the varying torus.

The componentwise exponential
\[
 \operatorname{Exp}_{N_{\mathrm{mesh}}}(a)_e
 =\exp(a_e/N_{\mathrm{mesh}})
\]
is only a local diffeomorphism at zero.  Its derivative is
\(N_{\mathrm{mesh}}^{-1}\operatorname{id}\), it is not globally injective,
and, when (DS_W(\mathbf1)=0\), at the critical identity configuration it
satisfies
\[
 D^2(S_W\circ\operatorname{Exp}_{N_{\mathrm{mesh}}})_0
 =N_{\mathrm{mesh}}^{-2}D^2(S_W)_{\mathbf1}.
\]

\Open\ The Gaussian space is an \(L^2\)-space on a noncompact linear
Coulomb slice with Lebesgue measure.  It is not the gauge-invariant subspace
of
\[
 L^2(\mathrm{SU}(2)^{E_{N_{\mathrm{mesh}}}},dU),
\]
the exponential supplies no global measure-preserving unitary map, and a
point mass at one link configuration is not in this Haar \(L^2\)-space.
No nonlinear compact-link Hamiltonian, time direction, reflection-positive
transfer matrix, continuum/infinite-volume/renormalized limit, or
Osterwalder--Schrader reconstruction is constructed.  Consequently the
finite Gaussian gap and small classical flux action do not imply any
physical four-dimensional quantum Yang--Mills spectral statement.  The
physical bridge remains open.  See
\cref{cusp-link-finite-spectrum,cusp-link-gaussian-theorem,%
cusp-link-exponential-local-only,cusp-link-physical-hilbert-space}.
\end{auditrecord}

\subsection{Cusp quotient and collapse}

The source defines an infinite toric threefold $Y$ from the height-one cones over the periodic $A_2$ triangulation and an action of the form
\[
\Psi_{\bar\lambda}(x,t)
=\bigl(c_{\bar\lambda}(t)t^{B_0\bar\lambda}x,t\bigr),
\qquad \bar\lambda\in\bar\Lambda.
\]

\begin{auditrecord}[C1: candidate equivariant cusp retraction]
\Verified\ The locally finite fan, smooth toric charts, action law, bounded
chart estimates, proper quotient, and fibre geometry are verified in
\cref{thm:candidate-cusp-geometry}.  The separate topological construction is
now proved in
\cref{lem:collapse-straightening,lem:positive-collapse,%
lem:phase-spread-collapse,prop:verified-cellular-collapse}: the positive
quotient is a compact manifold with collared central boundary, its retraction
lifts and spreads by phases, and straightening gives the equivariant diagram
\[
\begin{CD}
Y_{\le\eta} @>{r_Y}>> Y_0\\
@VV{/\bar\Lambda}V @VV{/\bar\Lambda}V\\
N'_0 @>{r}>> W.
\end{CD}
\]
with continuity across every toric stratum.  A tabulation of a separate finite
CW structure is unnecessary; the global equivariant retraction is the exact
claim used later.
\end{auditrecord}

\begin{auditrecord}[C2: integral collapse matrix]
\Verified\ With fixed primitive cellular orientations,
\cref{prop:verified-collapse-homology} proves the exact sequence
\[
0\longrightarrow
\operatorname{im}(\wedge^qM_0-I)
\longrightarrow\wedge^q\Lambda
\xrightarrow{c_*}H_q(W;\mathbb Z)
\longrightarrow0.
\]
It derives $c_*\colon\Lambda\to H_1(W)=\bar\Lambda$ as the quotient map and
proves
\[
c_*(\widehat\gamma\wedge\widehat u)=[\mathbb T_B],
\qquad
c_*(\widehat w\wedge\widehat\delta)=0,
\]
as well as the full mixed degree-two formula and the degree-three map in
explicitly oriented primitive bases.  Saturation is checked by the displayed
exterior-power image lattices, not inferred from ranks.  This result concerns
$c_*$; it does not identify the separate global Mayer--Vietoris matrices.
\end{auditrecord}

\begin{auditrecord}[C3: collapse equals specialization]
\Verified\ The intrinsic map and its saturation are proved in
\cref{thm:p0-integral-specialization}.  The functorial comparison is now
\cref{thm:collapse-equals-specialization}, which proves integrally that
\[
c^*=i_F^*r^*=\operatorname{sp}_{\mathrm{nc}}^q,
\qquad
\operatorname{im}c^*=(\wedge^qV)^{T_0},
\]
in every degree.  The proof realizes the nearby-cycle unit as
$i_F^*(i_W^*)^{-1}$ and uses the strong deformation retraction to identify the
inverse with $r^*$.  It therefore fixes the degree-three orientation and
excludes both a sign and a finite-index ambiguity without defining
specialization by the manuscript's answer.
\end{auditrecord}

\subsection{Based inclusions}

Choose one basepoint on the zero section and fixed whiskers through every collar.

\begin{auditrecord}[V1: primitive meridians and drift]
\Verified\ The actual coordinate gluings, one fixed zero-section base locus,
and explicit whiskers have no unrecorded conjugations and give
\[
\widehat g_j^{m_j}=t_{v_j},
\qquad
\widetilde a_0=t_\mu\widehat\rho_2^{-1}\widehat\rho_1^{-1}.
\]
The proof verifies that the clockwise meridian lifts to
$\widehat g_j$ rather than its inverse or a nonprimitive power, that the two
signs agree, that the zero-section whiskers add no hidden translation, and
that the tautological cusp disc gives $\mu=0$; see
\cref{prop:geometric-pi1}.  V1 is retained here as a resolved audit record.
\end{auditrecord}

\subsection{Mayer--Vietoris and Leray}

\begin{auditrecord}[MV1: geometric boundary cylinders]
\Verified\ The finite-band radial markings prove
$\alpha_q=(\pi_{1*},-\pi_{2*})$ as actual maps, with saturated degree-one
and degree-two matrices and top map $(3,-4)$.  For
$K_k=\ker(H_k(M)\to H_k(W))$, the toric disc and literal solid tori give
\[
K_1=\mathbb Z[\widetilde a_0],
\qquad
K_2=\langle\xi_{\widehat w},\xi_{\widehat\delta}\rangle,
\]
as primitive subgroups.  With the fixed orientation,
\[
\partial(j_*\xi_\nu)=(I-A_1^{-1})\nu=(I-A_2)\nu.
\]
The full cusp blocks reduce to $I_2\oplus[p]$ in degree one and
$I_4\oplus\operatorname{diag}(-p,1)$ in degree two; see
\cref{prop:actual-finite-band-maps,prop:cusp-mv-blocks}.
\end{auditrecord}

\begin{auditrecord}[L1: transgression comparison]
\Verified\ On the clockwise disc--pair-of-pants cover, the local discrepancy
for $12\gamma$ is
\[
\bigl(4\ell_1-b_1,\ 3\ell_2-b_2,\
-(12\ell_0-b_1-b_2)\bigr),
\]
whose signed sum is $-p$.  The transpose of the actual collar map produces
the negative of this same tuple.  The two primitive commuting squares in
\cref{thm:mv-leray-comparison} therefore identify the Mayer--Vietoris and
Leray obstructions, giving
\[
d_2(12\gamma)=-p\omega,\quad
d_2(2q)=-p\omega[\delta],\quad
d_2(2\gamma uw)=-p\omega[\gamma\delta].
\]
\end{auditrecord}

\subsection{Uniform non-torsion application}

\begin{auditrecord}[P10.8: one line bundle non-torsion on every fibre]
\Verified\ The candidate's period/countability route now gives a single
\(L\in\operatorname{Pic}(X)\) whose restriction is non-torsion on every
scheme-theoretic fibre and on both reduced finite fibres.  The workbench
proves
\[
R^1f_*\mathcal O_X\simeq\mathcal O\oplus\mathcal O(-1),
\qquad \operatorname{Pic}(X)\simeq\mathbb C,
\]
degree-one base change, and the exact orders
\(\operatorname{ord}(\mathcal O_X(S_j)|_{S_j})=m_j\).  The equation
\(c\tau+d\mu-b=0\) then yields the simultaneous countability conclusion.
The cusp extension is reconstructed from intrinsic nearby cycles and does
not use C1--C3 or MV1.  See
\cref{thm:structure-direct-images-audit,thm:uniform-nontorsion-audit}.
\end{auditrecord}

\subsection{All-Pic normalized cusp restriction}

\begin{auditrecord}[R1: Leray-edge restriction calculation]
\Verified\ For every $L\in\operatorname{Pic}(X)$, functoriality of the
Leray spectral sequence and proper base change identify the value of
$e(c_1(L))$ at $b$ with $c_1(L|_{X_b})$.  Since
\[
 \operatorname{im}e=E_\infty^{0,2}\subseteq\ker d_2^{0,2},
 \qquad d_2^{0,2}(2q)=\omega[\delta],
\]
the edge section is zero.  At $p_0$ this gives $c_1(L|_W)=0$, and
injectivity of $c_1$ on $\operatorname{Pic}(\mathrm{dP}_6)$ gives
$\eta^*i^*L\simeq\mathcal O$.  See
\cref{prop:leray-edge-restriction}.
\end{auditrecord}

\begin{auditrecord}[R2: boundary-Picard proof independent of global Leray]
\Verified\ For $E=c_1(L|_F)=2nq$, subtract the invariant vertical form on a
pair-of-pants core.  The resulting crossed-homomorphism boundary classes obey
\[
 \bar d_1+\bar d_2+\bar d_0=0,\qquad
 \bar d_1=4n[\delta],\quad
 \bar d_2=-3n[\delta],\quad
 \bar d_0=0.
\]
The finite values follow from the order-three/order-four norm maps and the
actual logarithmic translations; the cusp value follows from the explicit
zero-section sweep bounding a toric solid torus.  Thus $n=0$, and local
integral specialization gives
$\eta^*i^*L\simeq\mathcal O$.  See
\cref{lem:picard-crossed-defect,lem:picard-finite-norm,%
lem:picard-cusp-defect,prop:direct-boundary-picard}.
\end{auditrecord}

\begin{auditrecord}[R3: purely analytic restriction calculation]
\Verified\ The cusp-only version cannot work: the explicit divisor
$\mathcal D_Q$ in \cref{prop:local-cusp-picard-obstruction} descends to a
line bundle with generic class $-2q$ and nonzero normalized central degree
vector $(-12,2,0,-12,2,0)$.  Therefore the local map
$\operatorname{Pic}(N_0)\to\operatorname{Pic}(\mathrm{dP}_6)$ is nonzero.

\Verified\ The global calculation now incorporates both finite fillings.
For $c_1(L|_F)=kq$, logarithms of the actual $\mathcal O^*$ factors give
\[
 \bar d_1=2k[\delta],\quad
 \bar d_2=-\tfrac32k[\delta],\quad
 \bar d_0=0,\qquad
 \bar d_1+\bar d_2+\bar d_0=0,
\]
where the cusp equality follows from the argument principle in a Stein chart.
The finite sign certificate distinguishes every covering level.  Put
\[
 \widetilde J_j=\Delta_j\times\mathbb C^2,\qquad
 r_j:\widetilde J_j\longrightarrow J'_j,
\]
where the deck translations of the lattice quotient are $\tau_\lambda$.
The vector-cover automorphisms
\[
 Q_{j,\log}=\widetilde g_j^{\,v_j/m_j},\qquad
 Q_{j,\mathrm{lin}}=\widetilde g_j^{\,0}
\]
satisfy
\[
 Q_{j,\log}\tau_\lambda Q_{j,\log}^{-1}=\tau_{A_j\lambda},
 \quad (Q_{j,\log})^{m_j}=\tau_{v_j},\qquad
 Q_{j,\mathrm{lin}}\tau_\lambda Q_{j,\mathrm{lin}}^{-1}=\tau_{A_j\lambda},
 \quad (Q_{j,\mathrm{lin}})^{m_j}=1.
\]
 The positive sign is fixed by $m_j$ clockwise iterations.  Their quotient
 automorphisms $q_{j,\log},q_{j,\mathrm{lin}}$ on $J'_j$ have order dividing
 $m_j$ because $\tau_{v_j}$ is trivial after the lattice quotient, and order
 exactly $m_j$ because both cover $g_j$, whose order is $m_j$.
On $J_j'^*=J'_j|_{\Delta_j^*}$ the branch-independent translation $h_j$
obeys
\[
 h_jq_{j,\log}h_j^{-1}=q_{j,\mathrm{lin}}.
\]
With
\[
 p_{j,\log}:J_j'^*\to N_j^*,\qquad
 p_{j,\mathrm{lin}}:J_j'^*\to J^*,
\]
it induces
\[
 \bar h_j:N_j^*\xrightarrow{\ \sim\ }J^*,\qquad
 p_{j,\mathrm{lin}}\circ h_j=\bar h_j\circ p_{j,\log},\qquad
 p_{j,\log}^*(L|_{N_j^*})
 \simeq h_j^*p_{j,\mathrm{lin}}^*(L|_{J^*}).
\]
For the factor extension return to the filled vector cover.  If
\[
 P_{j,\log}:J'_j\to N_j,\qquad
 \pi_j^{\mathrm{vec}}:=P_{j,\log}\circ r_j:\widetilde J_j\to N_j,
\]
where \(P_{j,\log}\) is the full quotient and
\(P_{j,\log}|_{J_j^{\prime *}}=p_{j,\log}\), then its deck group is
\[
 \Gamma_j^{\log}
 =\langle\tau_\lambda,Q_{j,\log}:\lambda\in\Lambda\rangle,
\]
and $(\pi_j^{\mathrm{vec}})^*(L|_{N_j})$ is trivial on the contractible
Stein space $\widetilde J_j$.  Its factors define an integral central
extension, pushed out along $\mathbb Z\hookrightarrow\mathbb R$.  Over
$\Delta_j^*$ the line-bundle square identifies each fibre with the global
factor extension by a translation, so its commutator form is $E$.  This form
is integer-valued and locally constant on the connected disc $\Delta_j$,
hence remains $E$ on the filled cover.  Thus there are normalized vertical
lifts
\[
 U_\lambda^{(j)}U_\mu^{(j)}
 =\mathbf c\!\left(\frac12E(\lambda,\mu)\right)U_{\lambda+\mu}^{(j)}.
\]
For $\alpha\in\Gamma_j^{\log}$, define $A_\alpha,T_\alpha$ by
\[
 \alpha\tau_\lambda\alpha^{-1}=\tau_{A_\alpha\lambda},\qquad
 T_\alpha=(A_\alpha^{-1})^{\mathsf t},
\]
and define
$\widetilde d_j(\alpha)\in V_{\mathbb R}$ by
\[
 \widetilde G_\alpha
 U_{A_\alpha^{-1}\lambda}^{(j)}\widetilde G_\alpha^{-1}
 =\mathbf c\!\left(\left\langle\widetilde d_j(\alpha),\lambda\right\rangle\right)
 U_\lambda^{(j)},\qquad
 \widetilde d_j(\alpha\beta)
 =\widetilde d_j(\alpha)+T_\alpha\widetilde d_j(\beta).
\]
The additivity of the central exponent in $\lambda$ gives the unique
$\widetilde d_j(\alpha)$, and comparison of lifts of $\alpha\beta$ proves
the displayed crossed law.  To type the local-to-global comparison, put
\[
 \widetilde J_j^*=\Delta_j^*\times\mathbb C^2,
\]
\[
\begin{aligned}
 \Phi_{j,\log}
  &:=p_{j,\log}\circ r_j|_{\widetilde J_j^*},\\
 \Phi_{j,\mathrm{glob}}
  &:=p_{j,\mathrm{lin}}\circ h_j\circ r_j|_{\widetilde J_j^*}
   =\bar h_j\circ\Phi_{j,\log}.
\end{aligned}
\]
Both maps are $\Gamma_j^{\log}$-invariant.  Pulling the typed line-bundle
square to this common cover gives
\[
 \Theta_j:
 \Phi_{j,\log}^*(L|_{N_j^*})
 \xrightarrow{\ \sim\ }
 \Phi_{j,\mathrm{glob}}^*(L|_{J^*}),
\]
which induces an isomorphism of the pushed-out
$\Gamma_j^{\log}$-factor extensions and fixes every central
$\mathbf c(r)$.

Let $\eta_j$ be the total-space loop obtained by applying the zero section
to the based clockwise meridian $\rho_j$.  Choose
$x_j\in J_j^{\prime *}$ over $\eta_j(0)$ and put $y_j=h_j^{-1}(x_j)$.
The $p_{j,\mathrm{lin}}$-lift of $\eta_j$ starting at $x_j$ ends at
$q_{j,\mathrm{lin}}(x_j)$; its $h_j^{-1}$-transport is the
$p_{j,\log}$-lift of $\bar h_j^{-1}\eta_j$ starting at $y_j$ and ends at
$q_{j,\log}(y_j)$.  For a chosen vector lift $\widetilde y_j$, adjacent
logarithm branches satisfying
\[
 \ell'(w)=\ell(g_j^{-1}w)-\frac{2\pi i}{m_j}
\]
make the vector-cover endpoint exactly
$Q_{j,\log}(\widetilde y_j)$, with no inverse or lattice translate.
Consequently, for some $a_j\in V_{\mathbb R}$ and $b_j\in\mathbb R$,
\[
 \Theta_{j,*}(U_\lambda^{(j)})
 =\mathbf c(\langle a_j,\lambda\rangle)U_\lambda,\qquad
 \Theta_{j,*}(\widetilde G_{Q_{j,\log}})
 =\mathbf c(b_j)G_{\rho_j}.
\]
The central factor $\mathbf c(b_j)$ disappears under conjugation, so
\[
 \widetilde d_j(Q_{j,\log})=d_j+(T_j-I)a_j,\qquad
 \widetilde d_j(\tau_v)=\iota_vE,\qquad
 (\iota_vE)(\lambda)=E(v,\lambda).
\]
Consequently
\[
 N_{T_j}d_j
 =\widetilde d_j((Q_{j,\log})^{m_j})
 =\widetilde d_j(\tau_{v_j})=\iota_{v_j}E.
\]
Under $\ell\mapsto\ell+2\pi in$, the vector-cover lift changes by
$H_j\mapsto\tau_{nv_j}H_j$, while
\[
 \Theta'_{j,*}
 =\operatorname{Ad}(U_{nv_j})\circ\Theta_{j,*},\qquad
 a_j\longmapsto a_j+n\iota_{v_j}E.
\]
Since $A_jv_j=v_j$ and $E$ is $A_j$-invariant,
$T_j(\iota_{v_j}E)=\iota_{v_j}E$; hence both the correction
$d_j+(T_j-I)a_j$ and the norm identity are branch-independent.
This is the choice-independent typed extension (5.17-A14c)--(5.17-A15) in
\cref{prop:analytic-ah-boundary}.

Thus $k=0$.  For adjacent rays in a maximal cone
$\sigma=\{v,w,u\}$, the actual inclusions are
\[
 U_{\rho_v}=D(z_wz_u)\subset U_\sigma,\qquad
 U_{\rho_w}=D(z_vz_u)\subset U_\sigma,\qquad
 U_{\rho_v}\cap U_{\rho_w}=D(z_vz_wz_u)=U_0.
\]
The contractible Stein chart
\[
 S_\sigma=\{(z_v,z_w,z_u)\in\mathbb C^3:
 |z_vz_wz_u|<\epsilon\}
\]
has trivial Picard group and supplies the zero-horizontal-winding frame
comparison.  If
$\mu_\lambda^+:\Psi_\lambda^*\mathcal O(D_a)\to\mathcal O(D_a)$
is multiplication by $\chi^{m_\lambda}$, the descent tensor product uses
the inverse-dual map
\[
 \mu_\lambda^-=((\mu_\lambda^+)^{-1})^\vee:
 \Psi_\lambda^*\mathcal O(-D_a)\longrightarrow\mathcal O(-D_a),
 \qquad s\longmapsto\chi^{-m_\lambda}s.
\]
Put $M=L|_{N_0}$ and let
$E_M:=c_1(M|_{F_t})=c_1(L|_{F_t})$.  Also put
\[
 \Lambda_{\rm tor}:=\ker(M_0-I)=\operatorname{im}(M_0-I)
 =\mathbb Z\langle\widehat w,\widehat\delta\rangle,
 \qquad \overline\Lambda:=\Lambda/\Lambda_{\rm tor},
\]
\[
 B_0:\overline\Lambda\xrightarrow{\ \sim\ }\Lambda_{\rm tor},
 \qquad
 B_0\overline{\widehat\gamma}=-\widehat\delta,\qquad
 B_0\overline{\widehat u}=\widehat w.
\]
Its matrix is
$\left(\begin{smallmatrix}0&1\\-1&0\end{smallmatrix}\right)$ from
$(\overline{\widehat\gamma},\overline{\widehat u})$ to
$(\widehat w,\widehat\delta)$.
If $(d_1,d_2,d_3)$ are the independent degrees of
$\nu^*(M|_W)$, its six degrees are
$(d_1,d_2,d_3,d_1,d_2,d_3)$ and its toric support function has
polarization
\[
 B=\begin{pmatrix}d_2+d_3&d_3\\d_3&d_1+d_3\end{pmatrix},
 \qquad E_M((\lambda,0),(0,m))=-B(B_0\lambda,m)
 \quad(\lambda\in\overline\Lambda,\ m\in\Lambda_{\rm tor}).
\]
Because $k=0$ gives
$E_M=c_1(M|_{F_t})=c_1(L|_{F_t})=0$, the unimodularity of $B_0$ forces
$B=0$ and all six degrees to vanish, hence
$\eta^*i^*L\simeq\mathcal O$.  This third certificate avoids both global
Leray transgressions and local integral cusp specialization; see
\cref{prop:analytic-ah-boundary,lem:analytic-support-specialization,%
thm:analytic-all-picard}.
\end{auditrecord}

The intrinsic nearby-cycle/Leray route proves the recognition statement
\cref{cor:candidate-sphere}; the now-complete independent
collapse/Mayer--Vietoris route proves the same statement in
\cref{cor:candidate-mv-sphere}.  C1--C3, V1, MV1, and L1 are all resolved.

\subsection{Corrected CDP theorem}

\begin{auditrecord}[D1: strongest intact global criterion]
\Verified\ If one \(L_0\in\operatorname{Pic}(X)\) satisfies
\[
 H^2(X,T_X\otimes L_0)=0,
\]
then generic endpoint vanishing, semicontinuity, and HRR give
\[
 H^i(X,T_X\otimes M)=0\ (i\ne1),\qquad
 \chi(X,T_X\otimes M)=\tfrac12c_3(X)\le0.
\]
The relative vanishing \(R^2f_*(T_X\otimes L)=0\) for one \(L\) implies
this hypothesis by projection formula, Serre vanishing, and Leray.  For
reduced full fibres the exact replacement of the p.~692 step is injectivity
of
\[
 \delta_c:H^0(\operatorname{tors}\Omega^1_{W_c}\otimes A_c)
 \longrightarrow H^1(N^*_{W_c/X}\otimes A_c)
\]
together with the two endpoint vanishings.  See
\cref{thm:corrected-cdp-global,cor:corrected-cdp-relative,%
thm:corrected-cdp-fibrewise}.

This closes the corrected-replacement obligation.  It does not restore the
false theorem under its original hypotheses.  Because \(c_3(X)=2\), the
contrapositive of \cref{thm:corrected-cdp-global} gives
\(H^2(X,T_X\otimes L_0)\ne0\) for every \(L_0\); independently, the conductor
calculation gives \(R^2f_*(T_X\otimes L)\ne0\) for every \(L\), so the stronger
relative sufficient condition also fails.
\end{auditrecord}

\subsection{Exact Hodge-number calculation and replication boundary}

\begin{auditrecord}[H1: singular-fibre cotangent direct images]
\Verified\ The cyclic invariant lattices, toric logarithmic frame, and
residue sequence determine
\[
 R^0f_*\Omega_X^1=\mathcal O(-2),\qquad
 R^1f_*\Omega_X^1=\mathcal O\oplus\mathcal O(-1)^{\oplus2},\qquad
 R^2f_*\Omega_X^1=\mathcal O\oplus\mathcal O(-1).
\]
The logarithmic Kodaira--Spencer determinant has divisor $p_1+p_2$; its
degree-one exterior map is sheaf-surjective.  More exactly, on a finite
cover $t=s^m$, with $(m,a)=(3,2)$ or $(4,1)$ and
$D_{\mathrm{KS}}=(d_{rs})$, put
\[
 A=\frac{s}{m}d_{00},\qquad B=\frac{s^{1-a}}m d_{01},\qquad
 C=\frac{s^{1-a}}m d_{10},\qquad
 D_0=\frac{s^{m+1-2a}}m d_{11}.
\]
In the actual invariant bases,
\[
 [\kappa_0]=\begin{pmatrix}A&tB\\C&D_0\end{pmatrix},
 \qquad [\kappa_1]=(-C,A,-D_0,B).
\]
At the cusp, in the separately named logarithmic frames,
\[
 [\kappa_0]\equiv\begin{pmatrix}0&-1\\1&0\end{pmatrix}\pmod t,
 \qquad [\kappa_1]\equiv(-1,0,0,-1)\pmod t.
\]
The finite residue is
\[
 \operatorname{res}_{1,j}(xb_0+ys^ab_1)=m_jx[b_0],\qquad
 \operatorname{im}(\kappa_0\bmod t)
 =\ker\operatorname{res}_{1,j}=\mathbb C\,s^ab_1,
\]
so the two torsion stalks cancel against
$H^1(S_j,\mathcal O)$ under residue.  At the cusp put
\[
 q_1=(-1,0,1),\quad q_2=(0,-1,-1),\quad
 \alpha_i=\delta_0(q_i),\quad
 \theta=-\delta_1^{-1}(\alpha_1\smile\alpha_2).
\]
The standard \v Cech differential gives the indispensable negative sign
\[
 \delta_1\bigl(\rho_*(\alpha\otimes\eta_i)\bigr)
 =-\alpha\smile\alpha_i,
\]
and hence the exact row $(0,1,-1,0)$ in the ordered basis
$(\alpha_1\eta_1,\alpha_1\eta_2,\alpha_2\eta_1,\alpha_2\eta_2)$.

Finally the lift of $1$ through
$0\to\mathcal O(-1)\to f_*\mathcal E\to\mathcal O\to0$ is unique.
On the marked cover it is
\[
 h=(a(z),1)\binom{d\zeta_1}{d\zeta_2},\qquad a(z_2)=\frac12.
\]
The scalar must be evaluated in the genuine invariant frame
\[
 \omega_2=s\,u_2(s)(\gamma\wedge u)\ell,\qquad u_2(0)=1,
\]
not in $s(\gamma\wedge u)\ell$ away from $s=0$.  It is
\[
 \kappa_1(h\otimes\gamma)(p_2)
 =-\frac14\bigl(a(z_2)\dot\tau(z_2)+\dot\mu(z_2)\bigr)
 =-\frac14\dot\mu(z_2)\ne0,
\]
because $\dot\tau(z_2)=0$ while
$6\dot\mu(z_2)^2=K(z_2)\ne0$.  Therefore
\[
 h^{1,2}(X)=1,\qquad h^{1,1}(X)=2.
\]
See
\cref{prop:log-relative-direct-images,prop:log-ks-divisor,%
prop:absolute-cotangent-direct-images,thm:h12-scalar,cor:exact-hodge-diamond}.
\end{auditrecord}

\begin{auditrecord}[H2: independent Hodge replication]
\Open\ A further context-isolated AI derivation has independently recovered
the cyclic invariant lattices, logarithmic determinant divisor, residue
cancellation, order-four scalar, direct-image splittings, and
$h^{1,2}=1$, $h^{1,1}=2$.  Obtain an independent human specialist
replication of the same chain.  The internal proof and its numerical
conclusions are already \Verified; H2 is an external-replication boundary,
not a missing premise.
\end{auditrecord}

\begin{auditrecord}[H3: explicit representatives and products]
\Verified\ Global $\bar\partial$-closed representatives $G,Q$ have exact
marked-fibre period-coordinate classes and give the nine positive-degree
Bott--Chern generators
$\mathsf A,\mathsf B,\mathsf C,\overline{\mathsf C},
\mathsf D,\overline{\mathsf D},
 \mathsf E,\overline{\mathsf E},\mathsf T$.
On the marked fibre put
\[
 y=\operatorname{Im}\tau,\qquad m=\operatorname{Im}\mu,\qquad
 r=\operatorname{Im}\beta,\qquad
 \Delta=r-\frac{6m^2}{y}<0,
\]
\[
 \Gamma=\frac{i}{2\Delta}
 \left(d\bar\zeta_2-\frac{m}{y}d\bar\zeta_1\right),\qquad
 U=\frac{i}{2y\Delta}
 \left(r\,d\bar\zeta_1-6m\,d\bar\zeta_2\right).
\]
Then
\[
 \eta=u\wedge w+6\gamma\wedge\delta,
 \qquad \eta^2=12\gamma\wedge u\wedge w\wedge\delta,
\]
and its exact Dolbeault representative is
\[
 Q_F=-d\zeta_1\wedge U-6d\zeta_2\wedge\Gamma.
\]
The degree-one class is distinguished and the degree-two class is a chosen
coherent-Leray lift:
\[
 i_b^*\mathfrak g=[\Gamma]_{\bar\partial},\qquad
 i_b^*\mathfrak q=[Q_F]_{\bar\partial}.
\]
The choice of $\mathfrak q$ is unique only modulo
$\operatorname{im}d_1$, which restricts trivially to every smooth fibre;
accordingly an arbitrary smooth representative $Q$ need not restrict
coefficient-by-coefficient to $Q_F$.
The marked and complex orientations are opposite, and the exact pairings are
\[
 \int_{(F_b,\mathrm{or}_{\rm mark})}Q_F^2=12,
 \qquad
 \int_{(F_b,\mathrm{or}_{\mathbb C})}Q_F^2=-12,
 \qquad
 \int_{(F_b,\mathrm{or}_{\mathbb C})}
 Q_F\wedge\Gamma\wedge d\zeta_2=-1.
\]
With $\omega=c_1(\mathcal O_B(1))$, the integral Picard/Leray normalization
and Bott--Chern Chern formula are
\[
 \exp(12\mathfrak g)=f^*\mathcal O_B(1),\qquad
 c_{1,BC}(\exp(z\mathfrak g))=z\mathsf A+\bar z\,\mathsf B,
\]
so $c_{1,BC}(f^*\mathcal O_B(1))=12(\mathsf A+\mathsf B)$ in this
orientation.  Together with the multiplicative unit
$\mathbf1=[1]_{BC}$, the displayed positive-degree classes form bases in
the dimensions of
\cref{thm:exact-bott-chern-aeppli}, and their span $N$ satisfies
\[
 H^{*,*}_{BC}(X)=\mathbb C\mathbf1\oplus N,\qquad N^2=0.
\]
Indeed a $d$-closed function on compact connected $X$ is constant, and there
is no $\partial\bar\partial$ image of bidegree $(0,0)$, so
$H^{0,0}_{BC}(X)=\mathbb C[1]$.
See
\cref{prop:bcpr-generators-are-bases,thm:bcpr-square-zero-algebra}.  This
closes the class-and-product obligation.  A preferred literal global
coefficient atlas for the chosen smooth forms is not part of the claim.
The fibre pairings prove the required Dolbeault nonvanishing.  The vanishing
$\mathsf A\mathsf C=0$ uses Dolbeault--Serre duality together with the
independently proved injectivity
$H^{3,2}_{BC}(X)\to H^{3,2}_{\bar\partial}(X)$; the
Bott--Chern--Aeppli perfect pairing alone does not imply that product
vanishes.
\end{auditrecord}
